\section{Matched Comparison Runs}

After the xLSTM model-selection stage, the final controlled comparison used two
experiment launchers: \path{experiments/run_xlstm_tok14_batch16_sched_decay8000.sh}
for the selected xLSTM setup and
\path{experiments/run_transformer_tok14_gpt2_llama.sh} for the GPT-2-style and
LLaMA-style Transformer baselines. The corresponding run folders are
\path{experiments/xlstm-tok14-batch16-sched-decay8000} and
\path{experiments/transformer-tok14-batch16-sched-decay8000}. All three models
use the same train, validation, and test split files, tokenizer 14, a context
length of 1024 tokens, batch size 16, eight training epochs, and the same
learning-rate schedule from
\path{configs/train/transformer_tok14_batch16_sched_decay8000.yaml}. This gives
19,432 optimizer updates for each model.

The W\&B runs used for the plots are \verb|xlstm-tok14-batch16-sched-decay8000|
, \verb|tok_14_batch_16_gpt2| , and
\verb|tok_14_batch_16_llama| , all under the
\verb|duet-of-models| project. One logging detail matters for interpreting the
curves: the xLSTM trainer logs with the optimizer step as the W\&B step, whereas
the Transformer runs are logged through the Hugging Face trainer history, where
the W\&B row index is not the optimizer step. For the Transformer plots, the
optimizer step is therefore taken from the logged \verb|train/global_step|
field.

\subsection{Training and Test Metrics}

Figure~\ref{fig:matched-training-curves} compares training and validation
behavior for the three matched runs. All models improve rapidly during the first
few thousand updates, but the validation curves separate clearly. The xLSTM run
reaches the lowest validation loss at step 4000 and then remains relatively
stable. The GPT-2-style Transformer improves more slowly and reaches its best
validation loss at step 8000. The LLaMA-style Transformer reaches its best
validation loss at step 4000, after which its training loss continues to fall
while validation loss rises. This indicates stronger fitting capacity, but also
more overfitting under the fixed comparison schedule.

\begin{figure}[htbp]
	\centering
	\includegraphics[width=\textwidth]{figures/matched_training_curves.png}
	\caption{Training and validation loss, token accuracy, and top-5 token accuracy for the matched xLSTM, GPT-2-style, and LLaMA-style runs. The top row shows training metrics and the bottom row shows validation metrics, all plotted against optimizer step.}
	\label{fig:matched-training-curves}
\end{figure}

\begin{figure}[htbp]
	\centering
	\includegraphics[width=0.82\textwidth]{figures/matched_validation_loss.png}
	\caption{Validation-loss trajectories and selected checkpoints for the matched runs. The selected checkpoints are step 4000 for xLSTM, step 8000 for GPT-2-style Transformer, and step 4000 for LLaMA-style Transformer.}
	\label{fig:matched-validation-loss}
\end{figure}

Table~\ref{tab:matched-test-metrics} reports the held-out next-token metrics.
The xLSTM checkpoint is strongest across all reported test metrics: its test
loss is 1.3495, compared with 1.8555 for the LLaMA-style Transformer and 1.9657
for the GPT-2-style Transformer. The same ordering appears in perplexity and
token accuracy. The LLaMA-style model is the stronger Transformer baseline,
especially in token accuracy, but it does not close the gap to xLSTM under this
training setup. GPT-2 has similar top-5 accuracy to LLaMA, but its top-1
accuracy is substantially lower, meaning that it often places the reference
token among plausible alternatives without ranking it first.

\begin{table}[htbp]
	\centering
	\scriptsize
	\resizebox{\textwidth}{!}{%
		\begin{tabular}{lrrrrrl}
			\toprule
			Model                   & Best step & Val. loss       & Test loss       & Test ppl.       & Test acc.       & Selected checkpoint          \\
			\midrule
			xLSTM                   & 4000      & \textbf{1.4456} & \textbf{1.3495} & \textbf{3.8554} & \textbf{0.6148} & \path{best/checkpoint.pt}    \\
			LLaMA-style Transformer & 4000      & 1.8828          & 1.8555          & 6.3948          & 0.5150          & \path{llama/checkpoint-4000} \\
			GPT-2-style Transformer & 8000      & 2.0205          & 1.9657          & 7.1399          & 0.4306          & \path{gpt2/checkpoint-8000}  \\
			\bottomrule
		\end{tabular}
	}
	\caption{Held-out next-token metrics for the matched comparison runs. Lower loss and perplexity are better; higher token accuracy is better.}
	\label{tab:matched-test-metrics}
\end{table}

These results make the next-token comparison unambiguous. Under the matched
tokenization, context length, batch size, and training budget, xLSTM is the best
language model on the held-out test split. LLaMA is the better Transformer
baseline, but its validation curve suggests that the fixed schedule is not
optimal for it: by the end of training it has very low training loss but worse
validation loss than at the selected checkpoint. GPT-2 is more stable than
LLaMA, but plateaus at a higher validation and test loss.

\subsection{Generation Metrics}

The generation comparison uses the selected validation checkpoint from each run.
For the Transformer runs, Hugging Face Trainer reloads the best validation
checkpoint at the end of training and \path{trainer.save_model} writes that
model to the saved run root used by the generation launcher. The recorded best
checkpoints are \path{checkpoint-8000} for GPT-2-style Transformer and
\path{checkpoint-4000} for LLaMA-style Transformer. The xLSTM launcher uses the
explicit \path{best/checkpoint.pt} artifact. In all cases, generation uses the
shared evaluation configuration
\path{configs/eval/transformer_full_context_test.yaml}. The test split is
sampled in \verb|full_file_context| mode, the prompt is the maximum available
context up to 1024 tokens, and the continuation length is set to the remaining
reference length. Because \verb|generation_length| is set to \verb|reference|,
the configured \verb|max_generation_tokens| cap is not applied in these runs.
Sampling is stochastic with temperature 1.0, top-\(k=50\), and top-\(p=0.95\).

Generation time differs substantially between the models. Since the Transformer
runs did not store per-sample timing fields, their generation duration is
estimated from the modification-time span of the generated MIDI files. The same
timestamp estimate is reported for xLSTM, but the xLSTM run also contains
explicit timing telemetry for the resumed timed portion of generation. All three
models generated the same total continuation length, 3,798,456 tokens, so the
timestamp spans give a useful approximate throughput comparison.

\begin{table}[htbp]
	\centering
	\small
	\begin{tabular}{lrrl}
		\toprule
		Model                   & File span (h) & Approx. tokens/s & Explicit timing log      \\
		\midrule
		GPT-2-style Transformer & 5.50          & 191.9            & --                       \\
		LLaMA-style Transformer & 8.78          & 120.2            & --                       \\
		xLSTM                   & 44.81         & 23.5             & 36.20 h at 23.3 tokens/s \\
		\bottomrule
	\end{tabular}
	\caption{Approximate generation duration for the 200 generated MIDI files. File span is the time between the earliest and latest generated-MIDI modification timestamps in each collected output folder. The xLSTM run additionally logged explicit timing for 166 resumed samples.}
	\label{tab:matched-generation-timing}
\end{table}

The much longer xLSTM generation time is consistent with the cost of its
recurrent xLSTM blocks, especially the sLSTM blocks, during token-by-token
autoregressive decoding. However, this should be interpreted as a practical
runtime observation rather than a controlled architecture benchmark: the
timestamp spans also include implementation details, possible resume overhead,
I/O, and hardware/runtime differences between the runs.

Each model generated 200 continuations, and all generated outputs decoded and
scored successfully. The scoring pipeline compares the generated MIDI files
against their reference MIDI files with MusPy descriptors
\citep{dongMusPyToolkitSymbolic2020} and computes Frechet Music Distance (FMD)
\citep{retkowskiFrechetMusicDistance2025} over the generated and reference MIDI
sets. The FMD values are very close: GPT-2 has the lowest value at 201.91,
xLSTM follows at 202.54, and LLaMA has 205.64. Since all three values are within
a narrow range, the FMD result does not mirror the next-token ranking and should
be interpreted together with descriptor-level differences.

Figure~\ref{fig:matched-generation-metric-distributions} compares the generation
metrics with per-sample distributions for the MusPy descriptor differences.
FMD is shown as a separate set-level value because it is computed over the full
generated and reference MIDI collections rather than per individual sample.
The distributions show that the aggregate differences are not driven by a single
sample: GPT-2 has a wider upper tail for pitch-entropy and empty-beat-rate
differences, LLaMA has the highest median polyphony-rate difference, and xLSTM
keeps the lowest median empty-beat-rate and polyphony-rate differences while
having a higher pitch-range median.

\begin{figure}[htbp]
	\centering
	\includegraphics[width=\textwidth]{figures/matched_generation_metric_distributions.png}
	\caption{Distribution of the generation metrics for the matched runs. Violins show per-sample MusPy absolute differences, white boxes show the interquartile range and median, and the first panel shows FMD as a corpus-level metric. Lower values are better.}
	\label{fig:matched-generation-metric-distributions}
\end{figure}

\begin{table}[htbp]
	\centering
	\scriptsize
	\resizebox{\textwidth}{!}{%
		\begin{tabular}{lrrrrrrr}
			\toprule
			Model                   & Decode & Score & FMD             & Pitch ent. diff & Pitch range diff & Polyphony-rate diff & Empty-beat diff \\
			\midrule
			GPT-2-style Transformer & 1.00   & 1.00  & \textbf{201.91} & 1.2385          & \textbf{10.7350} & 0.3398              & 0.1313          \\
			LLaMA-style Transformer & 1.00   & 1.00  & 205.64          & \textbf{0.5764} & 11.3250          & 0.4027              & 0.0652          \\
			xLSTM                   & 1.00   & 1.00  & 202.54          & 0.6746          & 14.0800          & \textbf{0.2723}     & \textbf{0.0259} \\
			\bottomrule
		\end{tabular}
	}
	\caption{Generated-MIDI scoring summary. Decode and score columns are success rates. The remaining columns compare generated and reference MIDI sets, where lower is better.}
	\label{tab:matched-generation-metrics}
\end{table}

The generation metrics are more mixed than the next-token metrics. xLSTM has the
best rhythmic sparsity and polyphony-rate agreement among the three models: its
empty-beat-rate difference is 0.0259, compared with 0.0652 for LLaMA and 0.1313
for GPT-2, and its polyphony-rate difference is also the lowest. This suggests
that the xLSTM generations preserve some temporal density properties of the
reference set better than the Transformer baselines. However, xLSTM also has the
largest pitch-range and number-of-pitches differences, indicating that it tends
to cover a wider pitch space than the references.

The LLaMA-style Transformer gives the best pitch-distribution match among the
three models. It has the lowest pitch-entropy and pitch-class-entropy
differences, which means that its generated pitch usage is closer to the
reference distribution even though its FMD score is slightly worse and its
polyphony-rate difference is the highest. GPT-2 has the lowest FMD and the
smallest pitch-range difference, but it also has the largest pitch-entropy,
pitch-class-entropy, and empty-beat-rate differences. Therefore, its low FMD
should not be read as a complete win on generation quality; it is better on some
aggregate embedding statistics but weaker on several interpretable symbolic
descriptors.

Taken together, the two comparison stages lead to different conclusions. For
next-token language modelling, xLSTM is clearly the strongest model in this
matched setup. For generated MIDI, xLSTM remains competitive and is strongest on
rhythmic density and polyphony-rate agreement, while LLaMA is strongest on pitch
entropy and GPT-2 obtains the lowest FMD. The final model choice therefore
depends on the desired priority: choose xLSTM for the best held-out language
model and strong rhythmic agreement, choose LLaMA when pitch-distribution
similarity is the main criterion, and treat GPT-2's FMD result as useful but not
sufficient on its own because its descriptor-level deviations are larger.

\subsubsection{Hearing-Test Audio Renders}

For manual hearing checks, the generated MIDI folders\footnote{The generated MIDIs along with their WAVs can be found here \url{https://drive.google.com/drive/folders/1XJcGQ6y9gM19ZxW41jJVKB6jiRdwJzJ7?usp=sharing}} were partially rendered to
WAV with FluidSynth using the FluidR3 GM soundfont at 44.1 kHz. The renders were
stored next to the corresponding MIDI files with matching basenames. This
conversion was a practical preparation step for listening rather than a formal
blind listening study. It also exposed an important difference between the
models: many generated files were expensive to synthesize as uncapped WAV audio.
With a 60-second per-file render guard, GPT-2 completed 148 of 200 files, LLaMA
completed 105 of 200 files, and xLSTM completed 88 of 200 files. GPT-2 had the
highest completion count, but also the largest disk footprint: its completed WAV
payload occupied about 17 GB, with several files near or above 1 GB and one
timed-out retry expanding to 7.6 GB. LLaMA and xLSTM produced smaller completed
WAV sets, about 4.3 GB and 2.8 GB respectively, but both still left many files
unrendered under the timeout. For this reason, the hearing-test material should
be treated as a partial audio subset; rendering all 200 samples per model would
need duration caps or a compressed audio target.
